\section{Introduction}
\paragraph{What is metagenomics} Metagenomics is the application of modern genmic techniques to the study of communities of microbial organisms directly in their natural environment. This approach bypasses the need for isolating and cultivating individual species in the lab, which is particularly beneficial since many bacteria and viruses cannot survive in isolation.

In principle is possible to access 100\% of the genetic resources of an environment.

The applications are many:
\begin{itemize}
    \item Plant-microbe interaction for improving crop yield
    \item Monitoring and optimizing wastewater treatment
    \item Bioprospecting: hunting for novel enzymes and compounds
    \item Function and ecology of the oceans
    \item Development of functional feed
    \item Monitoring pathogens
    \item Monitoring polution
    \item Disease surveillance and treatment
\end{itemize}

The problem and our project
The comparison of metagenomics samples finds applications in precision medicine and environmental control. Given a set of metagenomics sample datasets, their all-against-all-comparison is a necessary step to derive a hierarchical clustering among the samples. However, it is also a very expensive step from the computational point of view.
In this project you are ask to investigates data mining techniques that avoid the computation of all-against-all-distances, such as, for example, PAM, CLARA, CLARANS or to independently think of an algorithmic approach, and test it on a metagenomics sample dataset.